\documentclass[11pt]{article}

\usepackage[a4paper,margin=2.4cm]{geometry}
\usepackage{mathptmx}
\usepackage[T1]{fontenc}
\usepackage[utf8]{inputenc}
\usepackage{amsmath,amssymb}
\usepackage{booktabs}
\usepackage{tabularx}
\usepackage{array}
\usepackage{microtype}
\usepackage{needspace}
\usepackage{float}
\usepackage{enumitem}
\usepackage{fancyhdr}
\usepackage[hidelinks]{hyperref}
\hypersetup{pdftitle={A Programmable Address Space for Neural Memory},
            pdfauthor={Y. Y. N. Li}}

\sloppy
\setlength{\emergencystretch}{3em}
\renewcommand{\arraystretch}{1.15}
\setlist[enumerate]{itemsep=2pt,topsep=3pt}
\setlist[itemize]{itemsep=2pt,topsep=3pt}
\newcolumntype{L}[1]{>{\raggedright\arraybackslash}p{#1}}

\pagestyle{fancy}
\fancyhf{}
\fancyhead[L]{\small\itshape A Programmable Address Space for Neural Memory}
\fancyfoot[C]{\small\thepage}
\renewcommand{\headrulewidth}{0.4pt}

\newenvironment{scopebox}{\begin{quote}\small\itshape}{\end{quote}}

\title{\bfseries A Programmable Address Space for Neural Memory}
\author{Y. Y. N. Li}
\date{}

\begin{document}
\maketitle
\thispagestyle{fancy}
\vspace{-1.2em}

\begin{abstract}
\noindent
Neural networks store knowledge in distributed parameters but do not normally expose an operational address
space whose contents can be selectively written, reassigned and transported while declared responses are
preserved. We assign the restricted response kernel a new operational role: not merely as a local source of
response-preserving directions, but as the carrier from which a finite memory-address structure can be
instantiated. We introduce Moving Fibre Intelligence, a framework that tests this role using ridge-regularized
near-kernel updates for operations within a response-constrained parameter set and local response retraction
for transport between declared response sets.

In separately frozen audits of pretrained GPT-2 using restricted LoRA-B parameter charts, writable endpoints
supported exact retrieval, overwrite and access through linguistic expressions excluded from within-run
writing updates. Two binary slots formed four compositional codes that were reused under cyclic
concept--code assignments. An eight-endpoint atlas spanning 28 pairwise relations was subsequently transported
to a distinct response set and returned while preserving exact and held-out-expression access. Mean relative
transport distortion was $1.45\times10^{-5}$, compared with $7.89\times10^{-5}$ for matched-scale
random-kernel controls, and direct retraction produced lower distortion in all 28 dependent pairwise
comparisons.

We then tested whether these resources could operate together in one shared parameter state. Four address
cells were jointly written, one was selectively overwritten, two exchanged their contents, and the resulting
state was moved to a second response set and returned. Every stage preserved exact and fresh
held-out-expression access, positive read margins and the declared response and KL budgets; where applicable,
unselected cells retained their declared values. Overwrite and swap also outperformed wrong-operation and
matched-norm random-update controls.

Within the tested GPT-2 LoRA charts, these results demonstrate an operational role for the restricted response
kernel as the local carrier of a finite compositional address geometry, with numerically near-kernel updates
supporting within-set operations and response retraction supporting local transport between sets. The procedures are externally
specified and gradient-based, and the evidence comes from separate single-seed protocols. Autonomous
instruction execution, unified multi-seed confirmation, capacity scaling, cross-model replication and global
transport remain open.
\end{abstract}

\section{Introduction}

Representation answers what internal state correlates with a concept; addressing asks whether a concept can
be assigned to and retrieved from explicit operational locations using reusable address codes. Attention, key--value caches,
sparse routing and external memory all implement addressing over activations or architecture. The question here concerns parameters: whether a response-constrained region of parameter space can
carry an explicit, auditable address code, whether its instantiated structure survives local transport, and
whether these resources can support selective operations on one shared parameter state.

Four traditions bear on this problem: distributed representation [2, 4], associative and key--value memory
[5, 12], model editing [8--10], and parameter symmetries and flat directions [1, 3, 6]. Within these traditions,
null directions have commonly appeared as redundancy or non-identifiability, and in some learning and editing
methods as constraints for preserving selected behaviour. What remains missing is their organization into an
explicit address geometry whose codes can be composed, reassigned and operated on in a shared parameter state.

The central mathematical shift of this paper is therefore not to discover that null directions can preserve
behaviour, but to test whether a restricted response kernel can serve as the local carrier of an explicit and
operable memory-address geometry. Concretely, we take
\begin{equation}
\mathcal{V}_\theta=\ker\bigl(DR_{A,S}(\theta)\big|_{T_\theta\mathcal{L}}\bigr)
\end{equation}
as that carrier in its ideal first-order form. Ridge-regularized near-kernel updates instantiate finite address
states within a declared response-constrained set, while local response retraction transports those
instantiated states between distinct declared sets. This kernel--retraction decomposition
provides the geometric basis for the compositional, remappable and shared-state operations tested below.

Operationally, we declare a finite response map on a fixed anchor set, take the parameters that keep that response within two
declared budgets, and write each concept--code pair to its own endpoint in that set; Section~2 gives the
definitions. Three construction steps then make the address atlas. Writing and overwriting
at a controlled endpoint establishes that these locations hold memories. Access through expressions excluded
from the writing updates establishes that a read is not a recital of a training string. Cyclic reassignment of
two-slot codes across concept identities establishes that a code is not the name of its concept. These three
are prerequisites, not parallel results.

These construction steps establish the resources required for an address atlas, but not yet their
transportability or shared-state operation. In a separate prospectively frozen transport audit, we therefore
ask whether an instantiated atlas preserves its readout and relational geometry between distinct
response-constrained parameter sets. Applying one local response-retraction rule separately to each endpoint of
an eight-endpoint cohort carried the cohort from one finite-response set to another while preserving every read
and, to within a declared aggregate gate, all 28 pairwise relations among the endpoints.

The final test is operational. We ask whether the preceding capabilities can be realized together in one
shared state by executing a frozen sequence of joint WRITE, selective OVERWRITE, SWAP, MOVE and RETURN.
OVERWRITE and SWAP are additionally audited against wrong-operation and matched-norm random-update controls.
This shared-state audit is the culminating experiment of the paper.

Together, these audits form a staged mechanism chain from reusable address construction, through local
transport, to selective shared-state operation within one restricted GPT-2 LoRA chart.

\section{Operational address geometry}

\subsection{Declared response map and admissible set}

Let $S$ denote the prospectively fixed selector of the declared response coordinates. For an anchor set
$A=\{x_1,\dots,x_m\}$, fixed before evaluation, the audited response map is
\begin{equation}
R_{A,S}(\theta)=\frac{1}{|A|}\sum_{x\in A} S f_\theta(x).
\end{equation}
Because the selector $S$ is fixed throughout each declared protocol, we suppress it in the notations
$\mathcal{F}_A$, $\mathcal{Y}_A$, $\mathcal{F}_s$ and $\mathcal{F}_t$. Response distances are measured in the
supremum norm on the selected coordinates, $d_R(y,y')=\lVert y-y'\rVert_\infty$, so a response budget bounds the
largest single-coordinate deviation rather than an aggregate.
For a declared target response $r$ and reference state $\theta_{\mathrm{ref}}$, the finite admissible set under
the response and collateral-KL budgets is
\begin{equation}
\mathcal{F}_A(r,\varepsilon)=\bigl\{\theta:\ d_R\bigl(R_{A,S}(\theta),r\bigr)\le\varepsilon_R,\quad
\overline{D}_{\mathrm{KL},A}(\theta_{\mathrm{ref}},\theta)\le\varepsilon_{\mathrm{KL}}\bigr\},
\end{equation}
where the collateral budget is the anchor-averaged divergence
\begin{equation}
\overline{D}_{\mathrm{KL},A}(\theta_{\mathrm{ref}},\theta)=\frac{1}{|A|}\sum_{x\in A}
D_{\mathrm{KL}}\bigl(p_{\theta_{\mathrm{ref}}}(\cdot\mid x)\,\Vert\,p_{\theta}(\cdot\mid x)\bigr).
\end{equation}
The KL reference is fixed prospectively within each protocol. In the staged shared-state audits it is the
initial pre-program state $\theta_{\mathrm{base}}$ and is not reset between operations; in the endpoint-wise
transport audit it is the corresponding pre-transport endpoint $\theta^{(s)}_{c,a}$.
The source chart is recovered by taking $r_s=R_{A,S}(\theta_{\mathrm{ref}})$; a target chart uses a separately
declared $r_t$, so that
\begin{equation}
\mathcal{F}_s:=\mathcal{F}_A(r_s,\varepsilon),\qquad \mathcal{F}_t:=\mathcal{F}_A(r_t,\varepsilon).
\end{equation}
These are finite budgeted response sets, not assumed to form a smooth global fibre bundle.

Let $\mathcal{L}\subset\mathbb{R}^p$ denote the restricted LoRA-B parameter chart, the only coordinates
updated in this work, and let $\mathcal{Y}_A$ denote the finite-dimensional response-coordinate space of the
declared anchor map. At $\theta\in\mathcal{L}$ the response differential is
$DR_{A,S}(\theta):T_\theta\mathcal{L}\to\mathcal{Y}_A$, and the local response kernel is
\begin{equation}
\mathcal{V}_\theta:=\ker\bigl(DR_{A,S}(\theta)\big|_{T_\theta\mathcal{L}}\bigr)
=\bigl\{v\in T_\theta\mathcal{L}: DR_{A,S}(\theta)v=0\bigr\},
\end{equation}
so that $R_{A,S}(\theta+\eta v)=R_{A,S}(\theta)+o(\eta)$ as $\eta\to0$ for $v\in\mathcal{V}_\theta$; no claim is made
about the kernel of the full GPT-2 parameter space.

In practice the kernel is realized numerically. With $J_\theta=DR_{A,S}(\theta)\big|_{T_\theta\mathcal{L}}$, the
ridge-regularized near-kernel operator used throughout is
\begin{equation}
P_{\theta,\lambda}=I-J_\theta^{\mathsf{T}}\bigl(J_\theta J_\theta^{\mathsf{T}}+\lambda I\bigr)^{-1}J_\theta,
\qquad \lambda=10^{-8}.
\end{equation}
We call $P_{\theta,\lambda}$ a numerical near-kernel operator rather than an exact projector, because
$\lambda>0$ generally implies $P_{\theta,\lambda}^2\neq P_{\theta,\lambda}$. Candidate tangent updates take the
form $v=P_{\theta,\lambda}u$ and lie near, rather than exactly in, $\mathcal{V}_\theta$; they are accepted only
after the declared finite response and KL gates are evaluated.

\subsection{Concept, code, endpoint, readout}

Let $\mathcal{C}$ denote the finite concept cohort and let
\begin{equation}
\mathsf{A}:=\{0,1\}^2
\end{equation}
be the four-element address-code set. An assignment is a declared map $\pi:\mathcal{C}\to\mathsf{A}$. Writing
concept $c$ under assignment $\pi$ produces an endpoint $\theta_{c,\pi(c)}$ with auditable readout
$\sigma(c,\pi(c))$, and the endpoint atlas instantiated by $\pi$ is
\begin{equation}
\mathcal{E}_\pi:=\bigl\{\theta_{c,\pi(c)}:c\in\mathcal{C}\bigr\}\subset\mathcal{F}_s .
\end{equation}
The load-bearing relation is therefore
\begin{equation}
(c,\pi(c))\ \longmapsto\ \theta_{c,\pi(c)}\ \longmapsto\ \sigma(c,\pi(c)).
\end{equation}
The same code may occur at several concept-specific endpoints, so an address code does not identify a unique
global parameter point. A slot is one coordinate of the code, and slots are audited separately.

\paragraph{Response kernel as a local carrier.}
Three levels should be kept apart: $\mathcal{V}_\theta$ is the ideal first-order local carrier,
$P_{\theta,\lambda}u$ is the numerically projected near-kernel candidate direction, and
$\operatorname{Retract}(\theta+\eta P_{\theta,\lambda}u)\in\mathcal{F}_A(r,\varepsilon)$ is the finite state
that the declared budgets actually admit. Accordingly, $\mathcal{V}_\theta$ is the local carrier of the tested
address operations; it is neither the discrete code set $\mathsf{A}$ nor the instantiated endpoint atlas
$\mathcal{E}_\pi$.

\subsection{Shared state and operation family}

A shared memory state $\Theta_m\in\mathcal{F}_s$ realizes a content map
\begin{equation}
m:\mathsf{A}\to\{0,1\}
\end{equation}
simultaneously within one parameter state. This object is distinct from the endpoint atlas $\mathcal{E}_\pi$,
whose elements are separately written concept--code endpoints. With $a,b\in\mathsf{A}$, the target state
relations audited in Section~5 are
\begin{align}
\operatorname{WRITE}(\Theta,m) &\mapsto \Theta_m, &
\operatorname{OVERWRITE}_a(\Theta_m,v) &\mapsto \Theta_{m[a\leftarrow v]},\\
\operatorname{SWAP}_{a,b}(\Theta_m) &\mapsto \Theta_{m^{a\leftrightarrow b}}, &
\operatorname{MOVE}_{s\to t}\bigl(\Theta^{(s)}_m\bigr) &\mapsto \Theta^{(t)}_m,\\
\operatorname{RETURN}_{t\to s}\bigl(\Theta^{(t)}_m\bigr) &\mapsto \widehat{\Theta}^{(s)}_m. & &
\end{align}
Here $m^{a\leftrightarrow b}$ exchanges the values assigned to cells $a$ and $b$ and leaves every other cell
unchanged, and $m[a\leftarrow v]$ sets cell $a$ to $v$ with the other cells untouched, while
$\Theta^{(s)}_m\in\mathcal{F}_s$ and $\Theta^{(t)}_m\in\mathcal{F}_t$ realize the same declared content map $m$
in different response-constrained sets. These are the target state relations of the experimental procedures,
not an algebra proved closed over arbitrary states, and not an instruction language parsed by the model.
Whether any of them can be executed while the declared budgets hold is the question of Section~5.

\subsection{Local response retraction and transport}

Candidate within-set updates are selected by backtracking and followed by ridge-regularized response
correction; a proposal with no admissible step is a zero step. The correction iterates the following
ridge-regularized minimum-norm update for the linearized response equation:
\begin{equation}
\delta\theta=-J_\theta^{\mathsf{T}}\bigl(J_\theta J_\theta^{\mathsf{T}}+\lambda I\bigr)^{-1}
\bigl(R_{A,S}(\theta)-r\bigr),\qquad \lambda=10^{-8},
\end{equation}
applied for a declared number of iterations and stopped early once the response residual falls below
$10^{-7}$. It is the regularized response-correcting counterpart associated with the same Jacobian: the
near-kernel operator suppresses first-order response change, whereas this update removes the remaining
linearized response residual. MOVE and RETURN apply these declared iterations directly and are
accepted only when their post hoc gates hold. Applied endpoint
by endpoint, that retraction carries endpoints between finite-response sets,
\begin{equation}
\mathcal{T}_{s\to t}:\mathcal{E}^{(s)}_\pi\longrightarrow\mathcal{E}^{(t)}_\pi,\qquad
\mathcal{T}_{s\to t}\bigl(\theta^{(s)}_{c,a}\bigr)=\theta^{(t)}_{c,a},
\end{equation}
where $\mathcal{E}^{(s)}_\pi\subset\mathcal{F}_s$ and $\mathcal{E}^{(t)}_\pi\subset\mathcal{F}_t$. The
notation names an experimentally specified endpoint-wise retraction map on a finite audited cohort, not a
global connection or parallel transport.

\subsection{Protocol provenance}

The evidence comes from separately frozen protocols rather than one continuous
or joint multi-seed experiment, and no statistic is pooled across stages. L1 is
a single-seed development record; the archived standalone L2 protocol lacks
its original numerical result; and the L3, L4 and L5 records are separate
prospective single-seed confirmation candidates. Held-out-expression claims
are drawn from the archived L3--L5 records, not reconstructed from the missing
L2 result. Protocol identifiers, seeds, hashes, complete gate definitions and
evidence locations are reported in the corresponding evidence directories and repository manifests. A
subsequent post-failure development smoke combined representative capabilities from all stages in one
continuous shared-state sequence and passed all declared smoke gates; because its configuration was adapted
after earlier failures on the same seed, it is reported only as joint-realizability development evidence and
does not replace the separately frozen stage-specific records.

\section{Constructing an address atlas}

\subsection{Writable, readable and overwritable endpoints}

A declared value was written at a controlled endpoint and read back exactly; a second value was then written at
the same location and read in place of the first; both endpoints satisfied the declared budgets. Random-kernel,
sign-reversed and no-move arms bound the alternative that any sufficiently small parameter change would produce
the same reads. Only endpoints meeting both budgets enter any positive claim.

\subsection{Access through excluded expressions}

Expressions excluded from the within-run gradient updates and from the checkpoint-selection rule read the same
written values. The positive held-out-expression claim reported here is supported by the archived
compositional-addressing and transport records, both of which prospectively evaluate expressions excluded from
their within-run writing updates; the standalone expression-access development protocol is retained as
historical provenance but is not used as an independently archived numerical result. Because the protocols were
produced through adaptive development on the same experimental family, this remains single-seed evidence rather
than an independently replicated generalization test.

\subsection{Assignment reuse and current-assignment readout}

The address-code set is $\mathsf{A}=\{0,1\}^2$, represented as $\{00,01,10,11\}$ and formed from two
separately audited slots, with a mixed-slot
minimum-margin gate so that a code using both slots must be read with positive margin rather than on average.
Let $\pi_j:c\mapsto a$ be the $j$-th cyclic assignment. Each assignment is written independently,
$(c,\pi_j(c))\mapsto\theta^{(j)}_{c,\pi_j(c)}$, so the four assignments occupy sixteen separately initialized
endpoints, and the audited property is
\begin{equation}
\operatorname{Read}\bigl(\theta^{(j)}_{c,\pi_j(c)}\bigr)=\pi_j(c) .
\end{equation}
The code rule is reused across concept identities and readout follows the current assignment.

\begin{center}\small
\begin{tabular}{clccc}
\toprule
$j$ & $\pi_j$ & Exact & Held-out & Mixed-slot minimum margin \\
\midrule
0 & dog 00; Tokyo 01; Japan 10; apple 11 & 8/8 & 32/32 & 1.9989 \\
1 & dog 01; Tokyo 10; Japan 11; apple 00 & 8/8 & 32/32 & 2.3205 \\
2 & dog 10; Tokyo 11; Japan 00; apple 01 & 8/8 & 32/32 & 1.6406 \\
3 & dog 11; Tokyo 00; Japan 01; apple 10 & 8/8 & 32/32 & 1.8391 \\
\bottomrule
\end{tabular}
\end{center}

Across the sixteen endpoints every exact and held-out read was correct, the minimum mixed-slot training margin
was 1.6406 against a gate of 0.5, and every endpoint stayed inside the declared response and KL budgets. The
result constructs a reusable address atlas, but each concept--code pair still occupies a separately initialized
endpoint. A training-only semantic router, correct on 79 of 80 outer-held-out concepts, assigned all four
bridge instances to their categories and hence, through a fixed category-to-address map, to their codes.
Per-view counts, the individual routing margins and the pair-only and shuffled-label control values are in the
archived records.

\begin{scopebox}
This is assignment reuse on independently initialized endpoints, not in-place rebinding of a single endpoint
and not rejection of a former binding. Shared-state storage of four cells, and selective operations on them,
are audited separately in Section~5. The archived record labels this run a prospective single-seed confirmation
candidate requiring multi-seed confirmation.
\end{scopebox}

\section{Transport across response fibres}

\subsection{Cohort and operation}

With $\mathcal{F}_s$ and $\mathcal{F}_t$ as declared in Section~2.1, the source cohort is the endpoint atlas
\begin{equation}
\mathcal{E}^{(s)}_\pi=\bigl\{\theta^{(s)}_{c,\pi(c)}:\ c\in\mathcal{C},\ \pi(c)\in\mathsf{A}\bigr\}\subset
\mathcal{F}_s,\qquad |\mathcal{C}|=8 .
\end{equation}
The cohort holds eight concept-specific endpoints over four repeated two-slot codes, two concepts per code: the
eight endpoints are not eight distinct addresses, and the code set is $\mathsf{A}$, not a three-slot cube. The
retraction map of Section~2.4 is applied endpoint by endpoint, so the operation is endpoint-level while the
audit is cohort-level.

\subsection{Declared audited quantities}

Three quantities are audited: the readout $I_{\mathrm{read}}=\operatorname{Read}(\theta_{c,a})$,
target-response eligibility
\begin{equation}
I_{\mathrm{elig}}(\theta;r)=\mathbf{1}\bigl[d_R\bigl(R_{A,S}(\theta),r\bigr)\le\varepsilon_R\ \wedge\
\overline{D}_{\mathrm{KL},A}\bigl(\theta^{(s)}_{c,a},\theta\bigr)\le\varepsilon_{\mathrm{KL}}\bigr],
\end{equation}
where the collateral budget of the transport stage is recorded relative to that endpoint's own pre-transport
state rather than to a global base state, and the
pairwise geometry $I_{\mathrm{pair}}=\{d(\theta_{c,a},\theta_{c',a'})\}$. The response value itself is
not an invariant: the target set is displaced from the source by a declared shift of 0.08, so what is preserved
is eligibility for the set the endpoints are carried to, not $R_{A,S}(\theta_{c,a})$. Geometric preservation is measured by the mean relative
distortion over the $|\mathcal{P}|=28$ unordered endpoint pairs,
\begin{equation}
D_{\mathrm{geom}}=\frac{1}{|\mathcal{P}|}\sum_{(u,v)\in\mathcal{P}}
\frac{\bigl|d_t(\theta^{(t)}_u,\theta^{(t)}_v)-d_s(\theta^{(s)}_u,\theta^{(s)}_v)\bigr|}
{d_s(\theta^{(s)}_u,\theta^{(s)}_v)+\epsilon},
\qquad u=(c,a),\quad v=(c',a').
\end{equation}

\subsection{Random-kernel perturbation control}

For each endpoint the control first adds a numerically projected response-kernel direction,
\begin{equation}
\theta^{\mathrm{ctrl}}=\operatorname{Retract}_{r_t}\bigl(\theta^{(s)}+z\bigr),\qquad
z=P_{\theta,\lambda}\xi,\qquad
\lVert z\rVert=\bigl\lVert\theta^{(t)}-\theta^{(s)}\bigr\rVert,
\end{equation}
and then applies the same retraction rule to reach the target response set. The control matches the norm of the
added kernel perturbation, not the final displacement, so the comparison asks whether direct retraction
preserves cohort geometry better than retraction preceded by a matched-scale random tangent perturbation.

\subsection{Result}

\begin{center}\small
\begin{tabular}{lc@{\qquad}lc}
\toprule
Cohort & Value & Geometry & Value \\
\midrule
Concept-specific endpoints & 8 & Transport distortion $D_{\mathrm{geom}}$ & $1.451\times10^{-5}$ \\
Unique two-slot codes & 4 & Control median distortion & $7.888\times10^{-5}$ \\
Endpoint-pair relations & 28 & Distortion ratio & 0.184 \\
Response displacement & 0.08 & Descriptive pairwise wins & 28/28 \\
Exact access after transport & preserved & Largest single-pair error & $3.97\times10^{-5}$ \\
Held-out access after transport & preserved & Round-trip distortion & $1.313\times10^{-7}$ \\
\bottomrule
\end{tabular}
\end{center}

Every exact and held-out read survived transport; the mean distortion over the 28 endpoint-pair relations met
the declared aggregate gate, with a largest single-pair relative error of $3.97\times10^{-5}$; and composing
transport with its reverse returned each endpoint to itself,
$\mathcal{T}_{t\to s}\circ\mathcal{T}_{s\to t}(\theta_{c,a})\approx\theta_{c,a}$, at a round-trip distortion of
$1.313\times10^{-7}$. Direct retraction produced lower distortion than the control on all 28 pairs. Transport
shows that an instantiated atlas can move without losing its declared structure; it does not yet show that
several address cells can coexist and be selectively manipulated in one state.

\begin{scopebox}
Two qualifications travel with these numbers. The declared gate is aggregate, not per-pair: it applies to
$D_{\mathrm{geom}}$ rather than to each of the 28 relations separately. And the 28 comparisons share the eight endpoints, so they are not independent
Bernoulli trials: the exact sign-test value of $p=3.73\times10^{-9}$ that they would imply under independence
was a declared gate of the frozen protocol but is reported here as a protocol diagnostic only, never as a
significance result, in line with the statistical note archived with the run. The symbol $\approx$ above is
defined by the declared response, read and pair-geometry gates, and closure over one source--target pair does
not imply path independence.
\end{scopebox}

\subsection{Alternative explanations and transport controls}

The result is not explained by independent endpoint refitting alone, because
the audit requires preservation of all pairwise relations in addition to
per-endpoint eligibility and readout. Nor is it explained by step size: a
matched-scale random kernel perturbation followed by the same retraction
produced approximately five times greater mean distortion. Round-trip recovery
to $1.313\times10^{-7}$ further excludes a strongly lossy projection on the
tested cohort. These controls distinguish local structure-preserving transport
from ordinary reassignment or unconstrained small movement, while remaining
limited to one source--target pair. The remaining question is operational: can
write, overwrite, rearrangement and transport be composed on one shared
address state?

\section{Programmable operations on a shared address state}

The preceding sections establish writable locations, expression-level access, reusable compositional codes and
locally transportable geometry. We now audit whether the target state relations defined in Section~2.3 can be
realized as selective operations within one shared parameter state while satisfying the declared response and
KL budgets. The audit is a prospectively frozen single-seed confirmation candidate at seed 82931, run with a
fresh seed and fresh held-out expressions after adaptive development of the protocol family.

\subsection{Program stages}

Four binary cells are written jointly into a single shared state, one cell is then selectively overwritten, and
two cells exchange their contents. WRITE, OVERWRITE and SWAP are response-constrained projected-gradient
updates on the restricted LoRA coordinates; the base model stays frozen. Every stage is scored on exact reads and on fresh held-out
expressions, and must leave the declared values of unselected cells unchanged.

\begin{center}\small
\begin{tabular}{lccccc}
\toprule
Stage & Declared bits & Cells changed & Unselected preserved & Response drift & Anchor KL \\
\midrule
WRITE & 0011 & --- & --- & $1.53\times10^{-5}$ & 0.00374 \\
OVERWRITE at 01 & 0111 & \{1\} & yes & 0 & 0.00351 \\
SWAP of 00 and 11 & 1110 & \{0,3\} & yes & $7.63\times10^{-6}$ & 0.00399 \\
\bottomrule
\end{tabular}
\end{center}

The observed bits matched the declared bits at every stage, exact and held-out accuracy were 1.0 throughout,
and the smallest held-out margin over the three stages was 0.6154, at the initial write.

\subsection{Operation controls}

OVERWRITE and SWAP are each accompanied by two controls: a wrong operation, which executes a different declared
state transition, and a matched-norm random-kernel update of the same update magnitude. On the requested-cell evaluation the
overwrite reached exact and held-out accuracy 1.0, against 0.5 for the wrong instruction and 0.75 for the
matched-norm random update at an update norm of 10.0006. The swap likewise reached 1.0, against 0.5 for both
its wrong-instruction and matched-random controls at an update norm of 13.2654. The requested operation is
therefore not reproduced by executing some other declared operation, nor by a displacement of the same size in
the response kernel.

\subsection{Moving the shared state}

The shared four-cell state was then carried to a second response set and returned. The parameter displacement
was 0.0978, the target-response residual $1.53\times10^{-5}$ and the return residual $1.53\times10^{-5}$. Exact
and held-out accuracy were 1.0 both after transport and after return, and the positive held-out margin gate
held at both ends. Thus all five externally specified operations were successfully executed in the declared
sequence on the shared four-cell state under the frozen access, selectivity, response and KL gates, and all
sixteen declared gates of the audit passed.

\begin{scopebox}
The operation names label externally executed experimental procedures, not instructions parsed by the model.
WRITE, OVERWRITE and SWAP are response-constrained projected-gradient updates; MOVE and RETURN are local
response retractions. The audit is a single-seed confirmation candidate: it does not establish multi-seed or
model-family confirmation, autonomous instruction execution, gradient-free or inference-time memory operations,
a general instruction-set architecture, arbitrary program composition, long-horizon error bounds, concept-level
capacity scaling, in-place semantic rebinding beyond the declared binary content operations, or global parallel
transport.
\end{scopebox}

\section{Discussion}

\subsection{A discrete address system inside a continuous chart}

The experiments identify the response kernel not only as a set of response-preserving null directions, but as a
local carrier on which addressable memory operations can be realized. They therefore distinguish the carrier
$\mathcal{V}_\theta$, the finite admissible set $\mathcal{F}_A(r,\varepsilon)$, the endpoint atlas
$\mathcal{E}_\pi$ and the shared state $\Theta_m$: these objects play related but non-interchangeable roles. The first structural result is that a
continuous constrained parameter chart supports a finite discrete code system whose slots can be audited
separately and recombined. Concepts are instantiated at endpoints
$\theta_{c,a}$, the code is read from separately audited slots, and every positive claim is an audit of
declared finite quantities rather than an inferred capacity of the model at large.

\subsection{From concept binding to code reuse}

The second result is separation: the same compositional codes can serve different concept identities, so the
tested code is not merely another name for one concept. Readout follows the declared assignment at each
separately written endpoint, and the mixed-slot margin gate rules out a writer that resolves one slot and
guesses the other.

\subsection{From geometric object to operational object}

The shared-state result changes the role of the geometry. In the construction stages the address structure is
built and read; under transport it is moved; in the shared state it becomes the state space of selective
operations. The progression is from location, to code, to transport, and finally to programmability, and each
stage consumes what the previous one produced.

This is also what separates the study from model editing. The object is not a single factual modification but
the selective execution and composition of several operations on one finite address state, and neither control
reproduced the requested state transition: the effect is tied to operation identity rather than displacement
magnitude alone.

\section{Limitations}

\begin{itemize}
\item Every stage uses one pretrained GPT-2 model, one restricted LoRA-B chart
and single-seed evidence; multi-seed and cross-model confirmation are absent.
\item L1--L4 use separately initialized concept--code endpoints. Shared-state
storage is tested only for four binary cells in Section~5, not for scaled
concept-level capacity or in-place semantic rebinding.
\item The operations are externally executed gradient-based procedures, not
model-parsed or inference-time instructions. Arbitrary programs, long-horizon
error bounds and a general instruction-set architecture are untested.
\item L4 uses four repeated two-slot codes over eight endpoints and applies
transport endpoint by endpoint between one source--target pair at one response
displacement; it does not define a global connection or parallel transport.
\item The random control matches the added kernel perturbation rather than the
final displacement, and the 28 pairwise comparisons share eight endpoints;
the declared geometry gate is aggregate rather than inferential.
\item Router bridge instances come from its fitting cohort, the
category-to-address map is hand-declared, and address information enters via a
controlled prompt interface rather than emerging autonomously.
\item The standalone L2 result archive is unavailable. No numerical result is
reconstructed from prose; held-out-expression claims rely on the publicly
archived later-stage records.
\end{itemize}

\section{Conclusion}

We have traced a staged path from writable parameter-space locations to a finite programmable address space.
Exact anchoring and held-out-expression access establish readable locations; compositional reassignment
separates address code from concept identity; local transport preserves the instantiated atlas across
response-constrained sets; and a shared-state audit combines these resources into selective WRITE, OVERWRITE,
SWAP, MOVE and RETURN operations. Within one prospectively evaluated GPT-2 LoRA seed, all declared access,
selectivity, response, KL and control gates of that audit passed. The result is a mechanism candidate for
programmable neural memory, not an autonomous instruction set or a general neural computer. Multi-seed
replication, cross-model validation, capacity scaling, longer programs and instruction-conditioned execution
remain open.

\section*{Data and code availability}
\addcontentsline{toc}{section}{Data and code availability}

\begingroup\sloppy
Code, frozen configurations, result records and claim boundaries are available
at \url{https://github.com/papasop/fibre-nav} under the archived L1--L5
evidence directories. This manuscript audits repository commit
\texttt{c243b0da8b8a05a2ebaffa6d310f8ca33b0416d7}.
\endgroup

\needspace{4\baselineskip}
\section*{References}
\addcontentsline{toc}{section}{References}

\small\begin{enumerate}[leftmargin=1.8em,itemsep=1pt,parsep=0pt]
\item S. Ainsworth, J. Hayase, and S. Srinivasa. Git Re-Basin: merging models modulo permutation symmetries.
\emph{International Conference on Learning Representations}, 2023.
\item T. Bricken et al. Towards monosemanticity: decomposing language models with dictionary learning.
\emph{Transformer Circuits Thread}, 2023.
\item L. Dinh, R. Pascanu, S. Bengio, and Y. Bengio. Sharp minima can generalize for deep nets.
\emph{International Conference on Machine Learning}, PMLR 70:1019--1028, 2017.
\item N. Elhage et al. Toy models of superposition. \emph{Transformer Circuits Thread}, 2022.
\item M. Geva, R. Schuster, J. Berant, and O. Levy. Transformer feed-forward layers are key-value memories.
\emph{Conference on Empirical Methods in Natural Language Processing}, 2021.
\item R. Hecht-Nielsen. On the algebraic structure of feedforward network weight spaces. In \emph{Advanced
Neural Computers}, pages 129--135. Elsevier, 1990.
\item E. J. Hu, Y. Shen, P. Wallis, Z. Allen-Zhu, Y. Li, S. Wang, L. Wang, and W. Chen. LoRA: low-rank
adaptation of large language models. \emph{International Conference on Learning Representations}, 2022.
\item K. Meng, D. Bau, A. Andonian, and Y. Belinkov. Locating and editing factual associations in GPT.
\emph{Advances in Neural Information Processing Systems}, 2022.
\item K. Meng, A. Sharma, A. Andonian, Y. Belinkov, and D. Bau. Mass editing memory in a transformer.
\emph{International Conference on Learning Representations}, 2023.
\item E. Mitchell, C. Lin, A. Bosselut, C. Finn, and C. D. Manning. Fast model editing at scale.
\emph{International Conference on Learning Representations}, 2022.
\item C. Olah et al. Zoom in: an introduction to circuits. \emph{Distill}, 2020.
\item S. Sukhbaatar, A. Szlam, J. Weston, and R. Fergus. End-to-end memory networks. \emph{Advances in Neural
Information Processing Systems}, 2015.
\end{enumerate}

\end{document}
